\documentclass[11pt,a4paper]{article}
\usepackage[margin=24mm]{geometry}
\usepackage[T1]{fontenc}
\usepackage{lmodern,microtype,amsmath,amssymb,amsthm,booktabs,array}
\usepackage{xcolor,enumitem,fancyhdr}
\usepackage[colorlinks=true,linkcolor=blue!45!black,urlcolor=blue!55!black]{hyperref}
\definecolor{ink}{HTML}{19374C}
\pagestyle{fancy}\fancyhf{}\fancyhead[L]{\small Counterfactual curation}\fancyhead[R]{\small Memory repair, version 1}\fancyfoot[C]{\thepage}
\setlength{\headheight}{14pt}\setlength{\parindent}{0pt}\setlength{\parskip}{6pt}
\setlist{nosep,leftmargin=1.5em}
\newtheorem{theorem}{Theorem}[section]
\newtheorem{proposition}[theorem]{Proposition}
\newtheorem{lemma}[theorem]{Lemma}
\title{\color{ink}\textbf{Eliminating Quadratic Coefficient Storage}\\[5pt]\Large Algorithms, proofs and local verification\\for counterfactual curation repair}
\author{ACL 2027 research project}
\date{4 October 2026}
\begin{document}\maketitle
\textbf{Outcome.} The quadratic expansion in the original ridge-summary
implementation has been removed. The stronger joint-span algorithm represents
each coefficient through its current Gram/cross-moment range. Its exact-rank
persistent numeric storage is $O(E_h(d+c))$, in place of $O(Kd^2+Kdc)$.
An implementation-safe envelope uses initial eligibility $E_0$, because
floating rank residue and sticky dense fallbacks can outlive exact rank loss.

\textbf{Implemented, not just proposed.} Two compact summaries, a reduced
ridge decoder, deterministic projector-based basis coordinates, streaming
checkpoints, and a stronger payload comparator were implemented and run here.
The matrix factorization primitives are classical; the paper contribution
must rest on their curation-aware guarantee and empirical consequences.

\textbf{Measured result.} On the fixed Civil engineering graph at $d=768$,
the original dense coefficient arrays used 234.483 MB. Joint-span arrays use
1.039 MB: a \textbf{225.6-fold reduction} (99.56\%). These are owned numerical
arrays, not process peak or total physical memory. A fair compact payload
baseline uses only 0.256 MB, so the new summary is not a universal memory winner.

\textbf{Evidence scope.} The same 100 natural Civil Comments records were
used, with their original labels. Curation is fixed 128-dimensional lexical
hashing at threshold 0.6, independent of learner dimension. No synthetic
documents or labels, semantic encoder claims, or external compute were used.
These are engineering results; full NLP utility and semantic validation remain
outside this local pilot.

\textbf{State scope.} The fast implementation repairs decoded coefficients
and heads. Canonical basis coordinates remove one historical ambiguity,
but sticky dense/compact mode and floating arithmetic prevent a claim of
whole-state equality to fresh construction. The stronger canonical-state
variant requires additional, explicitly charged re-encoding.
\newpage
\section*{Measured comparisons}
All cases share the same records, graph, horizon eight, regularization
$\lambda=0.01$, request generator and independent retained-data oracle.
The seven-output cases use all seven original Civil toxicity annotations;
they introduce no fabricated target. MB below means $10^6$ bytes.
\begin{center}\small
\begin{tabular}{rrrrrr}\toprule
$d$ & $c$ & Dense & Gram factor & Joint span & Compact payload\\\midrule
64 & 1 & 1.699 & 0.119 & 0.071 & 0.022 \\
128 & 1 & 6.641 & 0.271 & 0.172 & 0.043 \\
768 & 1 & 234.483 & 1.628 & 1.039 & 0.256 \\
128 & 7 & 7.249 & 0.879 & 0.183 & 0.047 \\
768 & 7 & 238.133 & 5.278 & 1.047 & 0.260 \\
\bottomrule\end{tabular}\end{center}
Counts in the two new summaries are exact Python integers, charged in their
separate Python live-size estimates. The table excludes those objects,
identifier dictionaries, graph/corpus input, allocator slack and temporary
workspace. The original dense state includes its count coordinate; that
small difference is not responsible for the reduction.

The signed-Gram ablation still stores $H_J$ densely. Its seven-output cost
exposes why the joint range, rather than Gram-only compression, is needed.
The compact payload comparator retains eligible FP32 features and FP64 labels,
stores no permanent Gram/kernel/head, and chooses primal or dual ridge by the
smaller dimension. Its persistent rows are a different storage contract.

\subsection*{Verification actually run}
\begin{itemize}
\item 544 release checkpoints across three scalar-output dimension panels
and two seven-output panels passed against an independently rescored oracle.
The joint summary's largest head difference was $1.11\times10^{-15}$; moment difference
was $4.00\times10^{-15}$.
\item The compact decoder was actually invoked at all 544 checkpoints. Its
largest oracle head difference was $1.17\times10^{-15}$.
\item Another 100 sequential releases deleted the entire natural corpus.
Joint-summary head difference was at most $2.06\times10^{-14}$; the final head was exactly
the prescribed zero. This is a stronger stress than the short request horizon.
\item Twelve atomic-rejection checks and three order/batching comparisons
passed. Twenty checkpoints per compressed method additionally compared every
coefficient to the dense implementation, streamed one key at a time.
\item Independent coordinate-permutation/sign audits used three actual
coefficients. Streamed checkpoint roundtrip reproduced coefficients exactly,
owned disjoint buffers, and supported an identical subsequent repair.
\end{itemize}
These are tolerance checks, not a certified roundoff or privacy proof. The
full coefficient reference is the dense implementation; head/moment checks
use a separate fresh selection-and-training oracle.

\newpage
\section*{Whole-process memory and the decision it supports}
Fresh processes separately constructed each method and serviced four releases.
They loaded the same natural inputs and graph; they did not run oracle/audit
reconstructions. Joint span used its reduced decoder. Each number is one
observed high-water mark on this host, not a replicated performance estimate.
\begin{center}\small
\begin{tabular}{lrr}\toprule
Method & Setup peak (MiB) & Final peak (MiB)\\\midrule
Dense coefficients & 117.78 & 356.71 \\
Signed Gram factors & 117.74 & 135.42 \\
Joint span + reduced decoder & 117.78 & 121.28 \\
Compact eligible payload & 117.77 & 118.39 \\
\bottomrule\end{tabular}\end{center}
The final peak includes the Python/numerical runtime, original inputs,
construction, requests and decoder workspace. Setup peaks must not be
subtracted and interpreted as exact allocation attribution. The 225.6-fold
array reduction is therefore not a 225.6-fold whole-process RAM claim.

\textbf{Recommended implementation.} Use joint-span coefficients for the
aggregate-state method. Keep dense coefficients as an audit reference and
signed-Gram factors as an explanatory ablation. Keep the compact payload
method as a mandatory baseline whenever eligible-row retention is allowed.
No speedup conclusion follows from these unreplicated runs.

\textbf{What is solved.} The summary no longer allocates one full ambient
Gram per coefficient. It is built directly from bounded per-key workspaces,
supports source/record substitution through the existing key algebra, and
can decode a compact current coefficient without an ambient $d\times d$
matrix. The source extension here is algebraic; this new numerical run
validates record deletion because the Civil preview lacks native source IDs.

\textbf{What is not proved.} This is not optimal physical storage, superiority
to payload retention, automatic full-state canonicality, numerical exactness,
privacy, or evidence of semantic NLP quality. A large corpus can still have
many eligible records and substantial key metadata. Construction preflight
caps coefficient arrays only; lifecycle caps must use the uniform envelope
and reserve repair, decode, serialization and external input workspaces.

\textbf{Reproducibility.} The companion package includes all three new
implementations, the exact design, per-checkpoint records, boundary results,
source hashes, independent audit and process measurements. Earlier execution
results remain intact. The code has no required network access during replay.
\newpage
% Standalone fragment; requires amsmath, amssymb, amsthm, hyperref and the
% theorem/lemma/proposition environments used by the main theory document.
% No document preamble. All guarantees below explicitly state their model.
\section{Removing the quadratic feature-dimension factor from persistent repair state}
\label{sec:memory-factor-addendum}

\paragraph{Scope and attribution.}
The dense implementation assigns a packed $d\times d$ Gram matrix to every
polynomial key. This is unnecessary for ridge statistics. The previous
linear-sketch lower bound permits arbitrary independent statistic coordinates;
it does not establish a quadratic-in-$d$ lower bound for rank-one ridge
contributions. The matrix primitives used below---thin QR, signed spectral
factorization, and low-rank update/downdate recompression---are classical, not
our proposed novelty. See Matthew Brand, \emph{Fast Low-Rank Modifications of
the Thin Singular Value Decomposition}, Linear Algebra and its Applications
415(1):20--30, 2006, DOI \href{https://doi.org/10.1016/j.laa.2005.07.021}{10.1016/j.laa.2005.07.021};
the author's institutional record is
\url{https://www.merl.com/publications/TR2006-059}.
The contribution requiring evaluation is their integration with the
finite-horizon coefficient state, its source-deletion semantics, and complete
memory accounting.

\subsection{Intrinsic coefficient rank and exact storage}
Fix the current retained collection, remaining horizon $h$, a frozen feature
map $z_v\in\mathbb R^d$, and frozen responses $y_v\in\mathbb R^c$. For record
deletion, a record's owner is itself. For disjoint-source deletion its owner
is its unique source $o(v)$, and $B(v)$ is the set of distinct blocker owners.
Discard records with $o(v)\in B(v)$ and records with $|B(v)|>h$; let $E_h$
denote the number of remaining eligible records. Each eligible record has
the two possible coefficient endpoints $B(v)$ and $B(v)\cup\{o(v)\}$, with
signs $+1$ and $-1$, respectively. The second endpoint is absent if its
degree exceeds $h$.

For a live key $J$, write its exact coefficient as
\[
 G_J=\sum_v a_{Jv}z_vz_v^\top,\qquad
 H_J=\sum_v a_{Jv}z_vy_v^\top,\qquad
 n_J=\sum_v a_{Jv},\qquad a_{Jv}\in\{-1,0,1\}.
\]
Let $m_J=|\{v:a_{Jv}\ne0\}|$ and let $K$ be the number of live keys.
Then $K\le2E_h$ and $\sum_Jm_J\le2E_h$ before or after exact zero-key
removal. Counts $n_J$ are signed integers; only the selected-set count
$n_\varnothing$ must be nonnegative.

\begin{lemma}[Intrinsic rank envelope]
For every key,
\[
 r_J:=\operatorname{rank}(G_J)\le m_J,
 \qquad
 \rho_J:=\operatorname{rank}([G_J\ H_J])\le m_J.
\]
Consequently $\sum_Jr_J\le2E_h$ and $\sum_J\rho_J\le2E_h$.
\end{lemma}
\begin{proof}
Every column of $G_J$ and $H_J$ lies in the span of the $m_J$ incident
feature vectors. Summing the ranks and using at most two endpoints per
eligible record proves the result. The argument allows signed coefficients
and repeated source owners.
\end{proof}

\paragraph{The first implemented representation: signed Gram factors.}
Represent $G_J=F_JS_JF_J^\top$, where $F_J\in\mathbb R^{d\times r_J}$ and
$S_J$ is diagonal with entries $+1$ or $-1$. Retain $H_J$ densely and
$n_J$ exactly. Alternatively retain $G_J$ in packed symmetric form, choosing
whichever representation occupies fewer numeric bytes. Writing
$q=d(d+1)/2$, a conservative word count is
\[
 \sum_J\bigl[\min\{(d+1)r_J,q\}+dc+1\bigr]
 \le 2E_h\bigl[d+dc+2\bigr].
\]
Thus persistent numeric state is $O(E_hd(c+1))$, rather than $O(Kd^2+Kdc)$.
This is an exact-arithmetic statement with rank recomputation after each
collision. It does not imply better memory than every payload baseline:
the constant factors, response dimension, precisions, and decoder workspace
still matter. A negative or indefinite coefficient is valid and requires
signed factors; replacing it by a positive-semidefinite approximation changes
the query polynomial.

\begin{theorem}[Joint-span representation]
There is an exact coefficient-only representation whose persistent numeric
state uses at most
\[
 \sum_J\left[d\rho_J+\frac{\rho_J(\rho_J+1)}2+
                    c\rho_J+1\right]
 \le 2E_h\left(\frac32d+c+\frac32\right)
\]
real/integer words. With key indices and alive public identifiers, logical
storage is $O(E_h(d+c+h)+U)$ words, where $U$ is the number of alive public
deletion units. This representation supports the same finite-horizon record
and disjoint-source deletion polynomial without retained feature rows,
response rows, or the original blocker graph.
\end{theorem}
\begin{proof}
Choose an orthonormal basis $U_J\in\mathbb R^{d\times\rho_J}$ of
$\operatorname{range}([G_J\ H_J])$, and store
\[
 A_J=U_J^\top G_JU_J,\quad B_J=U_J^\top H_J,\quad n_J.
\]
Because $G_J$ is symmetric, its row and column spaces lie in the same
subspace. Hence $G_J=U_JA_JU_J^\top$ and $H_J=U_JB_J$ exactly. The displayed
count stores $U_J$, the packed symmetric $A_J$, $B_J$, and the count.
Use $\rho_J\le d$, $\sum_J\rho_J\le2E_h$, and $K\le2E_h$ for the upper
bound. Each key has at most $h$ identifier incidences; its dictionaries and
degree/inverse indices therefore cost $O(hE_h+E_h)$ logical words. No
record-level decomposition of $G_J$ or $H_J$ is retained.
\end{proof}

\paragraph{Why the joint range is necessary.}
The range of $G_J$ alone is insufficient. Two signed contributions with the
same nonzero feature $z$, opposite signs, and distinct labels can give
$G_J=0$ but $H_J=z(y_1-y_2)^\top\ne0$. A scheme that stores cross moments
only in the Gram range silently loses information in this case. The
joint-span representation covers it without storing the original labels.

\subsection{Repairs, canonicality, and work}
The polynomial-key update is unchanged: substitute deleted variables by one,
merge colliding keys, and remove keys of degree above the decreased horizon.
For a joint-span collision, let $Q$ span $[U_1\ U_2]$, set
$P_i=Q^\top U_i$, and form
\[
 A=P_1A_1P_1^\top+P_2A_2P_2^\top,\qquad
 B=P_1B_1+P_2B_2,\qquad n=n_1+n_2.
\]
If $V$ spans $\operatorname{range}([A\ B])$, replace the state by
$(QV,V^\top AV,V^\top B,n)$ and discard its predecessor representations.
The signed-Gram variant instead obtains a thin QR factorization
$[F_1\ F_2]=QR$, diagonalizes
$R\operatorname{diag}(S_1,S_2)R^\top$, and absorbs square roots of absolute
eigenvalues into the new factor. Exact zero eigenvalues can be discarded;
negative eigenvalues cannot.

\begin{proposition}[Exact repair and state scope]
In exact arithmetic these operations decode to the canonical coefficient
map of the retained dataset at its remaining horizon. If each factorization
uses a deterministic convention depending only on its decoded current
coefficient, the encoded state is also a function of that current map and
the public membership/horizon state. It does not retain an original basis.
\end{proposition}
\begin{proof}
The factor operations preserve the sum of the decoded coefficients exactly.
Substitution and decreased-horizon truncation therefore commute with
decoding. The finite-horizon polynomial identity identifies the result with
fresh retained-data construction. A deterministic encoding of that map
preserves this state identity.
\end{proof}

For example, a canonical joint basis can be obtained by processing the
columns of $[G_J\ H_J]$ in fixed order, orthogonalizing each against earlier
accepted columns, skipping exactly zero residuals, and fixing normalization
signs. These columns can be generated from a small current factor, so an
entire dense ambient matrix need not persist. For spectral encoding, order
distinct eigenvalues, then canonically orient each eigenspace by applying
its projector to the coordinate vectors in fixed order and orthogonalizing.
An arbitrary QR or eigensolver output is not automatically this canonical
state, especially for repeated eigenvalues. The numerical implementation
must therefore claim coefficient/head fidelity, not bytewise canonical
factor equality.

For colliding widths $p,q$, set $s=p+q\le2d$. A streamed factor collision
uses $O(ds+s^2+dc)$ additional numeric workspace and, conservatively,
$O(ds^2+s^3+s^2c)$ arithmetic. The persistent-state bounds do not promise
the old dense algorithm's $O(d^2)$ cost per collision: repeated
recompression can be slower. If $C$ collisions occur, charge the sum of
these actual per-collision bounds, in addition to the unchanged key-index
work. Process collisions one at a time rather than allocating a dense
matrix for every key. A common dense ridge decoder separately uses
$O(d^2+dc)$ workspace and typically $O(d^3+d^2c)$ factor/solve work.
Construction must aggregate directly into factors or bounded per-key
workspaces; constructing the old $K$-matrix state and compressing afterward
does not cure its construction peak. Checkpoint export and array copies
must also be charged.

\subsection{What the floating-point implementation actually guarantees}
\paragraph{Exact rank and finite-precision rank are different.}
Mathematical cancellation may leave tiny nonzero FP64 eigenvalues, while
thresholding them would introduce an additional approximation. Retaining
every numerically nonzero eigenvalue avoids a user-selected truncation
tolerance, but does not make factorization exact or make stored width equal
to mathematical rank. QR, eigendecomposition, reconstruction, and summation
still introduce roundoff. The current-$E_h$ bounds above belong to the
exact-rank theorem, not automatically to software that keeps numerical
residue or leaves formerly dense entries dense.

\begin{proposition}[Implementation-safe initial-state envelope]
Consider the signed-Gram implementation with these explicit rules:
(i) a factor-factor collision returns at most the sum of the input widths;
(ii) dense storage is chosen only if its numeric bytes do not exceed the
candidate factor bytes; (iii) a collision involving a dense Gram remains
dense; and (iv) updates only rename, merge, or prune coefficient keys.
Then persistent Gram-array bytes, and also the total Gram/cross/count-array
bytes, never increase during a valid deletion sequence. For FP64 factors,
int8 signs, FP64 cross moments, and int64 counts, streamed rank-one
construction with $E_0$ initially eligible records gives the envelope
\[
 S_{\rm numeric}(t)\le S_{\rm numeric}(0)
       \le 2E_0(8d+1+8dc+8)\quad\text{bytes}.
\]
This bound excludes Python objects, key metadata, allocator slack,
serialization duplication, and transient repair/decoder workspace.
\end{proposition}
\begin{proof}
A factor of width $r$ occupies $8dr+r$ bytes. A factor-factor collision has
width at most the sum of its input widths, so its factor bytes are at most
their sum; choosing a smaller packed representation preserves this bound.
If an input is dense, retaining one packed matrix costs no more than the
existing dense input. Cross arrays and counts need one copy per resulting
key, and key count cannot increase. Each original eligible record inserts
at most two rank-one endpoint contributions, proving the initial bound.
No argument here depends on exact spectral cancellation. In particular,
the statement uses $E_0$, not the smaller current $E_h$.
\end{proof}

\paragraph{A distinct safe envelope for joint spans.}
Joint-span numeric bytes need not be monotone: the small core has a
quadratic term, so merging two bases can increase core bytes even when their
total width does not increase. For a joint implementation, assign each
initial endpoint one abstract ancestry token. Tokens are used only for the
analysis; they need not identify or retain a payload. Merge token sets on
collisions and discard them on pruning. Current entries have disjoint
ancestry and use at most $2E_0$ tokens altogether. A factor entry of width
$\rho$ uses at most its token count $a$, and $\rho\le d$. Its packed-core
storage therefore satisfies
\[
 d\rho+\rho(\rho+1)/2+c\rho+1
 \le (3d/2+c+1/2)a+1.
\]
If a dense entry is created only when its byte count does not exceed the
candidate factor byte count, charge it to that same token bound. Keeping it
dense on later collisions cannot exceed its earlier charge, and its token
set only grows. Hence the persistent numeric envelope is
$2E_0(3d/2+c+3/2)$ words with packed cores. If software stores the entire
$\rho\times\rho$ core instead, use the larger valid envelope
$2E_0(2d+c+1)$ words. This is a bound on all future numeric allocations in
the persistent state, not on per-collision temporary allocation. An initial
measured joint-state size is not by itself a safe future memory cap.

The numeric envelope is not a full physical-memory or privacy theorem.
Persistent factors can encode recoverable record information whenever the
underlying coefficient does; hiding rows syntactically is not a privacy
guarantee. No predecessor factors, old basis, or payload snapshots may be
retained under the stated current-summary contract. Allocator remnants,
backups, and previously released models remain outside that abstract-state
contract. Whole-process peak RSS, counted owned arrays, serialized bytes,
and common decoder workspace should be reported separately.

\paragraph{Limits of a single shared span.}
A global basis $U\in\mathbb R^{d\times r}$ shared by all keys gives storage
$O(dr+K(r^2+rc))$. It is useful when $r\ll d$, as can happen when a tiny
pilot has fewer examples than feature dimensions. When a realistic corpus
spans all $d$ features, $r=d$ and this does not remove the $Kd^2$ term.
Moreover a basis frozen from the original corpus can retain deleted-only
directions. The per-key rank envelope above neither assumes global low rank
nor preserves such an original basis.

\paragraph{Claims supported by this addition.}
The exact algorithm eliminates the unavoidable-looking quadratic feature
factor from persistent ridge-summary storage, with $O(E_h(d+c))$ words for
joint spans. The implemented signed-Gram variant has an auditable
$O(E_0d(c+1))$ numeric envelope even without numerical rank truncation.
Neither statement establishes universal superiority to an eligible-payload
baseline, smaller whole-process peak memory, faster repairs, bitwise
canonical floating-point state, or a privacy guarantee. Those distinctions
are part of the algorithm specification, not qualifications to omit from
the empirical comparison.

\subsection{Canonical basis orientation and compact decoding}
\paragraph{The implemented projector-pivot convention.}
For an orthonormal $U\in\mathbb R^{d\times\rho}$, initialize
$V=U^\top\in\mathbb R^{\rho\times d}$. Repeatedly select the first column
attaining the largest residual squared norm, normalize that column to
$o_j$, and subtract its orthogonal projection from every column of $V$.
After $\rho$ steps let $O=[o_1\ \cdots\ o_\rho]$. The reoriented state is
\[
 U_{\rm can}=UO,\qquad A_{\rm can}=O^\top AO,\qquad
 B_{\rm can}=O^\top B.
\]
This is pivoted Gram--Schmidt on projected coordinate vectors, performed
in small basis coordinates. It need not form the ambient projector
$UU^\top$. In exact arithmetic there is a positive pivot at every step
before $\rho$, since $U^\top$ has row rank $\rho$.

\begin{proposition}[Rotation invariance of the basis convention]
Suppose two orthonormal bases represent the same subspace and coefficient:
$\widehat U=UR$, $\widehat A=R^\top AR$, and
$\widehat B=R^\top B$, with $R$ orthogonal. The projector-pivot convention
produces identical $(U_{\rm can},A_{\rm can},B_{\rm can})$ in exact
arithmetic. It depends on the current subspace and coefficient, not on
the original construction basis.
\end{proposition}
\begin{proof}
Initially $\widehat V=R^\top V$. Orthogonal transformations preserve all
residual column norms, so both procedures select the same pivot, including
the first-index tie rule. Its normalized vector satisfies
$\widehat o_j=R^\top o_j$. Orthogonal projection subtraction preserves
$\widehat V=R^\top V$ at the next step. Induction gives
$\widehat O=R^\top O$. Consequently $\widehat U\widehat O=UO$, and
substitution into the two small-core transformations proves equality.
\end{proof}

This proposition is conditional on the represented subspace. If numerical
residue keeps extra basis directions, two repair histories can represent
different approximate subspaces. Even on the same mathematical subspace,
roundoff and near-tied pivots can produce different numerical choices.
Thus the implemented orientation is a deterministic convention, not a
proof of bytewise canonical floating-point state or exact numerical rank.
It also does not justify retaining a larger original basis indefinitely.

\paragraph{Basis orientation does not canonicalize representation mode.}
The fast implementation can keep a coefficient dense after an earlier
collision, even when a fresh construction would choose a compact factor.
Likewise, an initial contributor upper bound can affect a construction-mode
decision. These choices can depend on repair history even in exact
arithmetic. Canonical basis orientation removes basis rotations within a
chosen subspace; it does not remove this dense-versus-compact mode history.
The fast implementation consequently claims decoded coefficient/head
equivalence, not literal canonical-state equality. To instantiate the
stronger encoded-state theorem, the mode, width, and basis must all be
chosen from the actual current coefficient, for example by recompressing
dense entries on collisions or performing a canonical export. Such a step
can require $O(d^3)$ work per dense entry and its temporary memory must be
charged. It is not a free consequence of the pivot-orientation routine.

\begin{proposition}[Compact current-head decoder]
For a nonempty selected set with exact count $n>0$, suppose its current
moments are $G=UAU^\top\succeq0$ and $H=UB$, with $U^\top U=I_\rho$.
For averaged ridge regularization $\lambda>0$, the unique fitted head is
\[
 W^*=(G+\lambda nI_d)^{-1}H
     =U(A+\lambda nI_\rho)^{-1}B.
\]
Given the stored orthonormal joint representation, this requires
$O(\rho^3+\rho^2c+d\rho c)$ arithmetic and
$O(\rho^2+\rho c+dc)$ additional numeric workspace. There is no additional
$O(d\rho^2)$ basis-construction charge at decode time. If the state uses a
dense fallback, the corresponding solve has $O(d^2+dc)$ workspace and
$O(d^3+d^2c)$ arithmetic.
\end{proposition}
\begin{proof}
The right-hand side is in $\operatorname{range}(U)$. The orthogonal
complement has zero right-hand side and positive regularization, so its
solution component is zero. Write $W=UZ$ and multiply the normal equation
by $U^\top$ to obtain $(A+\lambda nI_\rho)Z=B$. A Cholesky factorization,
small solves, and multiplication by $U$ give the stated costs. The
positive-semidefinite assumption applies to the current selected Gram,
not to every signed polynomial coefficient.
\end{proof}

The cases $n=0$ and $\rho=0$ use the previously specified zero-head
convention where applicable; expressions divided by $n$ below require
$n>0$. In floating-point arithmetic, approximate orthogonality can make
the compact equation differ slightly from the ambient normal equation.
An ambient residual can still be evaluated without materializing $G$:
$UA(U^\top\widetilde W)+\lambda n\widetilde W-UB$.

\begin{proposition}[Head error from moment error and solve residual]
Let $(G_*,H_*,n)$ be the true retained moments, where $G_*\succeq0$ and
$n>0$ is exact. Let $(\widetilde G,\widetilde H,n)$ be represented
approximate moments and let $\widetilde W$ be any computed head. Suppose
\[
 \|\widetilde G-G_*\|_{\rm op}\le\varepsilon_G,
 \qquad \|\widetilde H-H_*\|_F\le\varepsilon_H,
 \qquad
 \|(\widetilde G+\lambda nI)\widetilde W-\widetilde H\|_F
       \le\rho_{\rm res}.
\]
Then, with $W_*=(G_*+\lambda nI)^{-1}H_*$,
\[
 \|\widetilde W-W_*\|_F
 \le\frac{\rho_{\rm res}+\varepsilon_H+
                      \varepsilon_G\|\widetilde W\|_F}{\lambda n}.
\]
\end{proposition}
\begin{proof}
The true normal-equation residual equals the represented residual plus
$(G_*-\widetilde G)\widetilde W+(\widetilde H-H_*)$.
Apply the triangle inequality and
$\|(G_*+\lambda nI)^{-1}\|_{\rm op}\le1/(\lambda n)$.
No positive-semidefinite assumption on $\widetilde G$ is needed for this
error bound, although an actual solver still must return a finite head.
\end{proof}

A small residual alone does not bound moment approximation error. The
displayed inequality becomes a numerical certificate only if all three
input bounds include rigorous floating-point error allowances. Measured
differences against a fresh oracle and ordinary floating-point residuals
are diagnostic validation, not such a certificate. The exact count
assumption is essential; a count error also perturbs the regularization
term and cannot be silently omitted.

\end{document}
